\documentclass{beamer}

\usepackage{cuhkbeamer}

\begin{document}

% ====================== Slide 1: TITLE SLIDE (REVISED) ======================
\begin{frame}
  \title{LAE Future Frontiers \& Final Review}
  \author{Chiu Yu KO}
  \institute{Low Altitude Economics}
  \date{Week 12}
  \titlepage
\end{frame}

% ====================== Slide 2: FORMAT & OBJECTIVES ======================
\begin{frame}[t]
\frametitle{Learning Objectives}
By the end of this class, you should be able to:
\begin{itemize}
\item \textbf{Explain} how adding nodes increases possible connections without guaranteeing demand or value.
\item \textbf{Evaluate} digital coordination using total system cost, reliability, oversight, and governance.
\item \textbf{Assess} possible Hong Kong roles in financing, leasing, insurance, professional services, and data-enabled assurance.
\item \textbf{Apply} the course viability test to integrated quantitative scenarios.
\item \textbf{Develop} recommendations that remain consistent with operational, regulatory, financial, and social constraints.
\end{itemize}
\end{frame}

% ====================== Slide 3: ICEBREAKER ======================
\begin{frame}[t]
\frametitle{The Next Horizon}
\textbf{The final week asks how an LAE project moves from a controlled trial to a repeatable operating system.}
\begin{itemize}
\item A larger network may create more useful connections, but also more coordination, repositioning, infrastructure, and regulatory requirements.
\item Growth should therefore follow evidence on demand, reliable capacity, unit economics, safety, and stakeholder impact.
\end{itemize}
\vspace{0.3cm}
\textbf{Quick poll (2 min):} Which part of the value chain is most likely to create durable value: aircraft, operations, infrastructure, software, finance, or professional services?
\end{frame}

% ====================== Slide 4: CONCEPT PRIMER 1 ======================
\begin{frame}[t]
\frametitle{Concept Primer: From Corridors to Networks}
A network links several origins, destinations, operators, and supporting facilities.
\begin{itemize}
\item More nodes create more possible connections, but only some will have sufficient demand and feasible operations.
\item Network value depends on frequency, reliability, access, transfer quality, compatible systems, and customer needs.
\item Cross-boundary services also require coordination concerning permissions, customs, immigration, quarantine, data, communications, liability, and emergency response.
\item Empty repositioning, unbalanced demand, congestion, and duplicated infrastructure can offset network benefits.
\end{itemize}
\textbf{Managerial lesson:} Count possible links first, then test which links are useful, serviceable, and commercially viable.
\end{frame}

% ====================== Slide 5: STANDALONE FORMULA ======================
\begin{frame}[t]
\frametitle{Applied Tool: Counting Possible Connections}
\begin{block}{Undirected Connections}
\[
L=\frac{N(N-1)}{2}
\]
\end{block}
\begin{itemize}
\item $N$ is the number of active nodes and $L$ is the number of possible pairs.
\item If direction matters, the number of possible origin-destination links is $N(N-1)$.
\item The formula counts possibilities, not scheduled routes, flights, demand, profit, or network value.
\item The number of links grows quadratically with $N$, not exponentially.
\item Each proposed link still requires demand, capacity, infrastructure, operating permission, and acceptable cost.
\end{itemize}
\end{frame}

% ====================== Slide 6: BACK-OF-ENVELOPE (NETWORKS) ======================
\begin{frame}[t,shrink=11]
\frametitle{Back-of-the-Envelope: Possible and Viable Links}
\textbf{Illustrative network expansion:}
\begin{columns}[T]
\begin{column}{0.48\textwidth}\begin{block}{Five Nodes}
\begin{itemize}
\item Possible node pairs: $5(4)/2=10$
\item Suppose six pairs pass demand and feasibility screening
\item \textbf{Candidate links: 6}
\end{itemize}\end{block}\end{column}
\begin{column}{0.48\textwidth}\begin{block}{Twenty Nodes}
\begin{itemize}
\item Possible node pairs: $20(19)/2=190$
\item Suppose 15\% pass initial screening
\item \textbf{Candidate links: about 29}
\end{itemize}\end{block}\end{column}
\end{columns}
\textit{Takeaway: More nodes increase the option set, but viable service grows only where demand, capacity, regulation, and operations align.}
\vfill{\tiny Illustrative screening rate; not a forecast for the GBA.}
\end{frame}

% ====================== Slide 7: CONCEPT PRIMER 2 ======================
\begin{frame}[t]
\frametitle{Concept Primer: Digital Coordination and Governance}
Digital systems can support planning, authorization, identification, traffic information, monitoring, and conflict management.
\begin{itemize}
\item Automation may reduce repetitive manual work and improve consistency at higher volumes.
\item Operators and service providers still require defined responsibilities, qualified personnel, cybersecurity, assurance, incident response, and regulatory oversight.
\item Human roles may change rather than disappear, particularly for exceptions, emergencies, system degradation, and accountability.
\item Algorithmic decisions should be traceable, testable, secure, and subject to appropriate review.
\end{itemize}
\textbf{Managerial lesson:} Evaluate the complete socio-technical system, not simply the cost of an algorithm.
\end{frame}

% ====================== NEW SLIDE 7.5: ALGORITHMIC GOVERNANCE COST SAVINGS ======================
\begin{frame}[t]
\frametitle{Applied Tool: Digital-Coordination Business Case}
\begin{block}{Annual Net Saving}
\[
\begin{aligned}
\text{Net Saving} = {} & \text{Cost of Current Process} \\
& {} - \text{Cost of Digital Process} - \text{Transition Cost}
\end{aligned}
\]
\end{block}
\begin{itemize}
\item Include software, communications, cybersecurity, certification, human oversight, maintenance, service-provider fees, and exception handling.
\item Savings depend on annual activity, reliability, integration complexity, and the amount of manual work actually displaced.
\item Marginal digital cost may be low, but it is not generally zero.
\item The preferred design must also meet safety, resilience, privacy, and accountability requirements.
\end{itemize}
\end{frame}

% ====================== NEW SLIDE 7.6: HK INSURANCE MODEL ======================
\begin{frame}[t,shrink=11]
\frametitle{Back-of-the-Envelope: Value of Better Risk Evidence}
\textbf{Illustrative assumptions:} Annual premium without project-specific evidence is HK\$12.5M. Evidence collection and risk controls cost HK\$3M.
\begin{columns}[T]
\begin{column}{0.48\textwidth}\begin{block}{No Demonstrated Improvement}
\begin{itemize}
\item Annual premium: HK\$12.5M
\item Evidence and control cost: HK\$0
\item \textbf{Total: HK\$12.5M}
\end{itemize}\end{block}\end{column}
\begin{column}{0.48\textwidth}\begin{block}{Improved Risk Case}
\begin{itemize}
\item Illustrative premium: HK\$9M
\item Evidence and control cost: HK\$3M
\item \textbf{Total: HK\$12M}
\item Illustrative first-year saving: HK\$0.5M
\end{itemize}\end{block}\end{column}
\end{columns}
\textit{Takeaway: Better evidence may support underwriting and risk management, but premium reductions are not automatic and must be evaluated against control costs.}
\vfill{\tiny Illustrative assumptions; not an insurance quotation or a measured Sandbox X effect.}
\end{frame}

% ====================== Slide 8: SESSION 2 TITLE ======================
\begin{frame}[c]
  \begin{center}
    \huge\color{cuhkpurple}\textbf{Session 2}\\
    \vspace{0.5cm}
    \Large HK's Services Edge \& Global Standard Setting
  \end{center}
\end{frame}

% ====================== Slide 9: CONCEPT PRIMER 3 ======================
\begin{frame}[t]
\frametitle{Concept Primer: Possible Hong Kong Service Roles}
A 2025 Our Hong Kong Foundation projection assigns a large share of potential 2050 LAE value added to supporting services.
\begin{itemize}
\item Possible roles include finance, leasing, insurance, legal and professional services, certification support, data assurance, cybersecurity, research, and training.
\item Hong Kong's actual advantage must be tested through capability, cost, regulation, customer demand, access to regional markets, and competition.
\item Manufacturing, operations, infrastructure, and services are complementary rather than mutually exclusive layers.
\item High value added does not imply high margin, low risk, or automatic market leadership.
\end{itemize}
\textbf{Managerial lesson:} Identify a specific service, customer problem, capability, revenue model, and defensible source of advantage.
\end{frame}

% ====================== Slide 10: STANDALONE FORMULA ======================
\begin{frame}[t]
\frametitle{Applied Tool: A Simple Leasing Appraisal}
\begin{block}{Annual Lessor Cash Flow}
\[
CF_L=\text{Lease Receipts}-\text{Financing Cost}-\text{Maintenance}-\text{Insurance}-\text{Expected Credit Loss}
\]
\end{block}
\begin{block}{Investment Value}
\[
NPV=-\text{Purchase Price}+\sum_{t=1}^{T}\frac{CF_{L,t}}{(1+r)^t}+\frac{\text{Residual Value}_T}{(1+r)^T}
\]
\end{block}
\begin{itemize}
\item Depreciation is an accounting allocation; residual value and cash expenditures belong in the cash-flow appraisal.
\item The lessor retains asset-value, financing, technology, maintenance, repossession, legal, and lessee-credit risks.
\item Lease receipts are contractual, not risk-free or guaranteed in economic terms.
\end{itemize}
\end{frame}

% ====================== Slide 11: BACK-OF-ENVELOPE (LEASING) ======================
\begin{frame}[t,shrink=12]
\frametitle{Back-of-the-Envelope: Risk Does Not Disappear}
\textbf{Illustrative monthly scenario:} Lease payment HK\$300k.
\begin{columns}[T]
\begin{column}{0.48\textwidth}\begin{block}{Operator}
\begin{itemize}
\item Revenue after disruption: HK\$250k
\item Lease payment: HK\$300k
\item Other operating costs omitted
\item \textbf{Shortfall before other costs: HK\$50k}
\item Bears demand and operating risk
\end{itemize}\end{block}\end{column}
\begin{column}{0.48\textwidth}\begin{block}{Lessor}
\begin{itemize}
\item Contractual lease receipt: HK\$300k
\item Financing and asset-management cost: HK\$150k
\item Cash flow before credit loss: HK\$150k
\item \textbf{If the lessee defaults, receipt and asset value are at risk}
\end{itemize}\end{block}\end{column}
\end{columns}
\textit{Takeaway: Leasing reallocates risks between parties; it does not create a risk-free margin.}
\vfill{\tiny Illustrative assumptions; excludes maintenance, insurance, tax, residual value, and recovery cost.}
\end{frame}

% ====================== Slide 12: GLOBAL STANDARDS ======================
\begin{frame}[t]
\frametitle{From Local Evidence to Transferable Capability}
Controlled trials can generate evidence on aircraft performance, communications, navigation, weather, emergency response, infrastructure, simultaneous operations, and user feedback.
\begin{itemize}
\item Evidence may support local policy development, operating procedures, assurance services, and commercial refinement.
\item Transfer to another jurisdiction is not automatic because rules, environments, aircraft, routes, and acceptance conditions differ.
\item A credible exportable capability may include testing methods, data governance, assurance, training, technical integration, or professional advice.
\item Claims of a global ``gold standard'' require external validation, reproducibility, interoperability, and recognition by relevant authorities.
\end{itemize}
\textbf{Managerial lesson:} Export validated methods and capabilities, not an assumption that one local rulebook will be universally adopted.
\end{frame}

% ====================== Slide 13: SESSION 3 TITLE ======================
\begin{frame}[c]
  \begin{center}
    \huge\color{cuhkpurple}\textbf{Session 3}\\
    \vspace{0.5cm}
    \Large Quantitative Applications \& Final Exam Prep
  \end{center}
\end{frame}

% ====================== Slide 14: EXAM STRATEGY ======================
\begin{frame}[t]
\frametitle{The Final Exam: A Five-Step Method}
For an integrated LAE scenario:
\begin{enumerate}
\item \textbf{Define the decision}: What is being chosen, for whom, over what period, and compared with what alternative?
\item \textbf{Find served activity}: Compare customer demand with reliable capacity.
\item \textbf{Calculate viability}: Use price, variable cost, served activity, fixed cost, and any specified charges or payments.
\item \textbf{Map each constraint once}: Show whether it changes demand, capacity, price, cost, timing, or risk.
\item \textbf{Recommend conditionally}: Address safety, regulation, sustainability, distribution, uncertainty, and the evidence needed to proceed.
\end{enumerate}
\textit{The aim is not to force every course tool into every answer. Use only the tools that clarify the decision.}
\end{frame}

% ====================== Slide 15: STANDALONE FORMULA ======================
\begin{frame}[t]
\frametitle{Applied Tool: The Viability Test}
\begin{block}{Step 1: How Much Can We Actually Serve?}
\[
\text{Served Activity}
=
\min\{\text{Customer Demand},\text{Reliable Capacity}\}
\]
\end{block}
\begin{block}{Step 2: Does the Service Cover Its Costs?}
\[
\text{Operating Profit}
=
(\text{Price}-\text{Variable Cost})
\times\text{Served Activity}
-\text{Fixed Cost}
\]
\end{block}
\begin{itemize}
\item Add any specified tax, fee, subsidy, or committed payment in the correct place.
\item Recalculate served activity when a shock changes demand or capacity.
\item Then discuss regulation, safety, sustainability, and distributional effects.
\end{itemize}
\end{frame}

% ====================== NEW SLIDE 15.5: QUANTITATIVE TOOLKIT RECAP ======================
\begin{frame}[t]
\frametitle{Quantitative Toolkit Recap}
Choose the tool that matches the question:
\begin{itemize}
\item \textbf{Demand}: generalized cost, choice share, elasticity, and market funnel.
\item \textbf{Capacity}: reliable flights, seats or payload, operating days, and bottlenecks.
\item \textbf{Economics}: contribution, breakeven activity, unit cost, cash flow, and NPV.
\item \textbf{Uncertainty}: scenarios, sensitivity analysis, expected value, probability of loss, and liquidity need.
\item \textbf{Allocation}: queueing, reservations, priority rules, congestion charges, and auctions.
\item \textbf{Impact}: life-cycle emissions, noise exposure, privacy controls, stakeholder distribution, and KPIs.
\end{itemize}
\textit{Start with the decision and units. Do not begin with an all-purpose ``super-formula.''}
\end{frame}

% ====================== Slide 16: BACK-OF-ENVELOPE (EXAM PREP) ======================
\begin{frame}[t,shrink=12]
\frametitle{Back-of-the-Envelope: A Capacity and Cost Shock}
\textbf{Base data:} Demand = 200 flights/day; reliable capacity after disruption = 160; price = HK\$500; variable cost = HK\$100/flight; fixed cost = HK\$40,000/day; specified noise charge = HK\$50/flight.
\begin{columns}[T]
\begin{column}{0.48\textwidth}\begin{block}{Step 1: Served Activity}
\[
\min\{200,160\}=160
\]
\textbf{Step 2: Contribution}
\[
(500-100-50)\times160
=\text{HK\$56,000}
\]
\end{block}\end{column}
\begin{column}{0.48\textwidth}\begin{block}{Step 3: Operating Profit}
\[
56{,}000-40{,}000
=\textbf{HK\$16,000/day}
\]
\begin{itemize}
\item The service remains profitable under the stated assumptions.
\item Compare mitigation cost with the expected reduction in the charge and other benefits.
\end{itemize}\end{block}\end{column}
\end{columns}
\textit{Do not recommend visual masking as a noise remedy. Privacy and noise require different controls.}
\vfill{\tiny Illustrative exam assumptions; not a current Hong Kong noise charge or capacity rule.}
\end{frame}

% ====================== Slide 17: POLICY CONSISTENT JUDGMENT ======================
\begin{frame}[t]
\frametitle{The Ultimate Goal: Consistent Judgment}
A strong answer combines correct calculation with a recommendation that respects the operating context.
\begin{itemize}
\item Keep units consistent and distinguish demand, capacity, served activity, flights, seats, and passenger-trips.
\item Do not apply the same constraint twice through separate ``haircuts.''
\item Do not assume that Sandbox X bypasses legal or safety requirements; it supports controlled testing and evidence collection.
\item Do not assume a market mechanism is appropriate where safety or predefined public-service priority governs access.
\item State which assumptions are provided, which are sourced, and which are analytical scenarios.
\item Conclude with the condition under which the project should proceed, pause, redesign, or stop.
\end{itemize}
\end{frame}

% ====================== Slide 18: CONCLUSION ======================
\begin{frame}[t]
\frametitle{Conclusion: From Analysis to Decision}
The Low Altitude Economy remains an emerging field in which technology, regulation, business models, and public expectations continue to develop.
\begin{itemize}
\item Good analysis separates evidence from assumptions.
\item Good modelling connects demand, reliable capacity, cost, finance, risk, and impact without double counting.
\item Good strategy uses staged commitments and evidence-based decision gates.
\item Good policy considers safety, efficiency, innovation, competition, accessibility, privacy, environmental impact, and distribution.
\end{itemize}
\vspace{0.3cm}
\begin{center}
\textit{The objective is not to predict the future with certainty, but to make better decisions as evidence develops.}
\end{center}
\end{frame}

% ====================== Slide 19: THANK YOU ======================
\begin{frame}[t]
  \frametitle{Thank You + Final Q\&A}
  
  \begin{center}
      \Large \textbf{End of Course}
  \end{center}
  
  \vspace{1cm}
  \textbf{Final Reminders:}
  \begin{itemize}
      \item \textit{Review all past reading bundles and calculation frameworks.}
      \item \textit{Final Exam logistics and timing.}
  \end{itemize}
  
  \vspace{0.5cm}
  \begin{center}
      \textbf{Questions for the Final Exam?} \\
      \small (Final 15 minutes reserved for extended Q\&A)
  \end{center}
\end{frame}

\end{document}